\chapter{Linear Combination of Unitaries and Block-Encoding Arithmetic}
% Solutions to the exercises of \cref{ch:lec07}.

\begin{solution}[Exercise~\ref{ex:lec07-1}]
Applying the three factors in turn,
\[
  (\mathrm{PREP}\otimes\id)\ket{0^a}\ket\psi
  =\frac{1}{\sqrt{\norm{c}_1}}\sum_j\sqrt{c_j}\ket{j}\ket\psi
  \xrightarrow{\mathrm{SEL}}
  \frac{1}{\sqrt{\norm{c}_1}}\sum_j\sqrt{c_j}\ket{j}U_j\ket\psi.
\]
Projecting the control register onto $\bra{0^a}$ after $\mathrm{PREP}\dg$ uses
$\bra{0^a}\mathrm{PREP}\dg=(\mathrm{PREP}\ket{0^a})\dg
=\frac{1}{\sqrt{\norm{c}_1}}\sum_{j'}\sqrt{c_{j'}}\bra{j'}$, so with
$\braket{j'|j}=\delta_{jj'}$,
\[
  (\bra{0^a}\otimes\id)W(\ket{0^a}\otimes\id)\ket\psi
  =\frac{1}{\norm{c}_1}\sum_{j,j'}\sqrt{c_{j'}c_j}\,\delta_{jj'}U_j\ket\psi
  =\frac{1}{\norm{c}_1}\sum_j c_j U_j\ket\psi=\frac{T}{\norm{c}_1}\ket\psi.
\]
For a general input flag $\ket{k}$, the first step gives
$(\mathrm{PREP}\otimes\id)\ket{k}\ket\psi=\sum_j (\mathrm{PREP})_{jk}\ket{j}\ket\psi$,
then $\mathrm{SEL}$ inserts $U_j$, and projecting the output onto $\bra{0^a}$ with
$\bra{0^a}\mathrm{PREP}\dg=\frac{1}{\sqrt{\norm{c}_1}}\sum_j\sqrt{c_j}\bra{j}$ leaves
\[
  (\bra{0^a}\otimes\id)W(\ket{k}\otimes\id)
  =\frac{1}{\sqrt{\norm{c}_1}}\sum_j\sqrt{c_j}\,(\mathrm{PREP})_{jk}\,U_j.
\]
Only for $k=0^a$, where $(\mathrm{PREP})_{j0}=\sqrt{c_j}/\sqrt{\norm{c}_1}$, does this
collapse to $T/\norm{c}_1$; the encoded operator depends solely on the $\ket{0^a}$
column of $\mathrm{PREP}$, as claimed.
\end{solution}

\begin{solution}[Exercise~\ref{ex:lec07-2}]
With $\mathrm{PREP}\ket{0^a}=\frac{1}{\norm{\beta}_2}\sum_j\beta_j\ket{j}$ and
$\widetilde{\mathrm{PREP}}\ket{0^a}=\frac{1}{\norm{\gamma}_2}\sum_j\gamma_j\ket{j}$,
the same computation as in Exercise~\ref{ex:lec07-1} yields
\[
  (\bra{0^a}\otimes\id)(\widetilde{\mathrm{PREP}}\dg\otimes\id)\mathrm{SEL}
  (\mathrm{PREP}\otimes\id)(\ket{0^a}\otimes\id)
  =\frac{1}{\norm{\beta}_2\norm{\gamma}_2}\sum_j\gamma_j\beta_j U_j
  =\frac{T}{\norm{\beta}_2\norm{\gamma}_2},
\]
since $\beta_j\gamma_j=c_j$. Thus the subnormalization is
$\alpha=\norm{\beta}_2\norm{\gamma}_2$. By Cauchy--Schwarz,
\[
  \norm{c}_1=\sum_j\abs{\beta_j\gamma_j}
  \le\Bigl(\sum_j\beta_j^2\Bigr)^{1/2}\Bigl(\sum_j\gamma_j^2\Bigr)^{1/2}
  =\norm{\beta}_2\norm{\gamma}_2,
\]
with equality iff $(\abs{\beta_j})$ and $(\abs{\gamma_j})$ are proportional, i.e.\
$\abs{\beta_j}=\abs{\gamma_j}$. Combined with $\beta_j\gamma_j=c_j\ge0$ this forces
$\beta_j=\gamma_j=\sqrt{c_j}$, at which point $\widetilde{\mathrm{PREP}}=
\mathrm{PREP}$ and $\alpha=\norm{c}_1$. Hence among all such circuits the symmetric
choice attains the minimum subnormalization $\norm{c}_1$.
\end{solution}

\begin{solution}[Exercise~\ref{ex:lec07-3}]
Here the coefficients are $c=(1,1,2)$ for $(U_0,U_1,U_2)=(\Xgate,\Ygate,\Zgate)$, so
$\norm{c}_1=1+1+2=4$. Three terms need a control register of $a=2$ qubits ($K=4$),
with label $j=3$ unused ($c_3=0$). The prepare state is
\[
  \mathrm{PREP}\ket{00}=\frac{1}{\sqrt4}\bigl(\sqrt1\ket{0}+\sqrt1\ket{1}
  +\sqrt2\ket{2}+0\cdot\ket{3}\bigr)
  =\tfrac12\ket{0}+\tfrac12\ket{1}+\tfrac{1}{\sqrt2}\ket{2},
\]
(a valid unit vector: $\tfrac14+\tfrac14+\tfrac12=1$), and the select oracle is
\[
  \mathrm{SEL}=\proj{0}\otimes\Xgate+\proj{1}\otimes\Ygate
  +\proj{2}\otimes\Zgate+\proj{3}\otimes\id
\]
(the $\ket3$ branch is irrelevant since its amplitude is zero). Then
$W=(\mathrm{PREP}\dg\otimes\id)\mathrm{SEL}(\mathrm{PREP}\otimes\id)
\in\mathrm{BE}_{4,2}(T)$, i.e.\ $(\alpha,a)=(4,2)$. To check the overshoot, $T=\Xgate
+\Ygate+2\Zgate$ is Hermitian with $T^2=(\Xgate+\Ygate+2\Zgate)^2$; using
$\Xgate^2=\Ygate^2=\Zgate^2=\id$ and the anticommutation
$\{\Xgate,\Ygate\}=\{\Ygate,\Zgate\}=\{\Zgate,\Xgate\}=0$, all cross terms cancel and
$T^2=(1+1+4)\id=6\id$, so $\norm{T}=\sqrt6\approx2.449$. Thus $\alpha=\norm{c}_1=4$
overshoots $\norm{T}=\sqrt6$; the gap $4/\sqrt6\approx1.63$ is the price of the
LCU decomposition (a single-unitary encoding would give $\alpha=\sqrt6$).
\end{solution}

\begin{solution}[Exercise~\ref{ex:lec07-4}]
\emph{Equal factors.} With $U_A,U_B\in\mathrm{BE}_{\alpha,a}(\cdot)$ sharing the flag
and one control qubit ($m=1$), take $\mathrm{PREP}\ket0=\tfrac{1}{\sqrt2}(\ket0+
\ket1)=\Hgate\ket0$ and $\mathrm{SEL}=\proj0\otimes U_A+\proj1\otimes U_B$. By
\cref{prop:lec07-sum} with $\alpha_A=\alpha_B=1$ (take $A,B$ with $\norm{\cdot}\le1$),
$W=(\Hgate\otimes\id)\mathrm{SEL}(\Hgate\otimes\id)$ encodes
$\tfrac12 A+\tfrac12 B$; to encode $C=A+B$ instead, use weights $c_0=c_1=1$, i.e.\
$\gamma=\alpha_A+\alpha_B=2$ and the same $\Hgate$ prepare, giving
$W\in\mathrm{BE}_{2,\,a+1}(A+B)$.

\emph{Unequal factors.} For $U_A\in\mathrm{BE}_{\alpha_A,a}(A)$,
$U_B\in\mathrm{BE}_{\alpha_B,a}(B)$, set $\gamma=\alpha_A+\alpha_B$ and
\[
  \mathrm{PREP}\ket0=\frac{1}{\sqrt{\gamma}}\bigl(\sqrt{\alpha_A}\ket0
  +\sqrt{\alpha_B}\ket1\bigr).
\]
With $B_A=(\bra{0^a}\otimes\id)U_A(\ket{0^a}\otimes\id)=A/\alpha_A$ and likewise
$B_B=B/\alpha_B$, the projected block is
\[
  (\bra0\bra{0^a}\otimes\id)W(\ket0\ket{0^a}\otimes\id)
  =\frac{\alpha_A}{\gamma}B_A+\frac{\alpha_B}{\gamma}B_B
  =\frac{1}{\gamma}(A+B),
\]
so $W\in\mathrm{BE}_{\alpha_A+\alpha_B,\,a+1}(A+B)$: subnormalization
$\alpha_A+\alpha_B$ and $a+1$ ancilla qubits (the shared flag plus one control).
\end{solution}

\begin{solution}[Exercise~\ref{ex:lec07-5}]
Write $U_B\ket0\ket\psi=\ket0 B\ket\psi+\ket1\ket\chi$ (the $\ket1$ branch collects
the unflagged part). Then
\[
  U_A U_B\ket0\ket\psi
  =U_A\bigl(\ket0 B\ket\psi+\ket1\ket\chi\bigr)
  =\bigl(\ket0 A B\ket\psi+\ket1\ket{\chi'_1}\bigr)
   +U_A\ket1\ket\chi,
\]
and $U_A\ket1\ket\chi$ generally has a nonzero $\ket0$-component, say
$\ket0\,R\ket\chi+\ket1\ket{\chi'_2}$ with $R=(\bra0\otimes\id)U_A(\ket1\otimes\id)$.
Hence
\[
  (\bra0\otimes\id)U_AU_B(\ket0\otimes\id)\ket\psi=AB\ket\psi+R\ket\chi,
\]
which differs from $AB\ket\psi$ by the cross term $R\ket\chi$. Concretely take
$A=B=\tfrac12\Xgate$, block-encoded (e.g.\ by $U_A=U_B=\Hgate$-conjugated controlled
gates); since $\norm{A}=\tfrac12<1$ the branch $\ket\chi\ne0$, and $R\ne0$ for a
generic such encoding, so $(\bra0\otimes\id)U_AU_B(\ket0\otimes\id)\ne\tfrac14\Xgate^2
=\tfrac14\id=AB$. Placing the encodings on \emph{disjoint} flag qubits removes the
mixing: by \cref{prop:lec07-prod},
\[
  (\bra{00}\otimes\id)(U_A\otimes\id)(\id\otimes U_B)(\ket{00}\otimes\id)
  =\frac{A}{1}\cdot\frac{B}{1}=AB,
\]
because projecting register $a$ onto $\ket0$ commutes through $\id\otimes U_B$ and
projecting register $b$ onto $\ket0$ commutes through $U_A\otimes\id$, so the two
blocks factor cleanly.
\end{solution}

\begin{solution}[Exercise~\ref{ex:lec07-6}]
Iterating \cref{prop:lec07-prod} over $L$ identical factors
$U_A\in\mathrm{BE}_{\alpha,a}(A)$ on disjoint flag registers gives
$\prod_{i=1}^{L}(U_A\text{ on register }i)\in\mathrm{BE}_{\alpha^L,\,aL}(A^L)$; the
compression gadget (\cref{rem:lec07-compress}) reduces the ancilla count to
$a+\lceil\log_2(L+1)\rceil$ while keeping the same subnormalization $\alpha^L$. For a
nonnormal example take
\[
  A=\begin{pmatrix}0&\tfrac32\\ 0&0\end{pmatrix},\qquad
  \rho(A)=0\ (\text{both eigenvalues }0),\quad \norm{A}=\tfrac32>1,
\]
which is nilpotent: $A^2=0$, so $\norm{A^L}=0$ for $L\ge2$, bounded (indeed zero).
Yet any block encoding must have $\alpha\ge\norm{A}=\tfrac32$, so encoding $A^L$ via
the product rule carries subnormalization $\alpha^L\ge(3/2)^L\to\infty$. The
mismatch---bounded $\norm{A^L}$ but exploding $\alpha^L$---arises because the product
rule multiplies the \emph{worst-case} single-step subnormalizations and is blind to
the cancellation that makes the true power small; it is the block-encoding analogue of
the gap between $\norm{A}$ and $\rho(A)$ for nonnormal $A$, and motivates the
dedicated differential-equation algorithms of \cref{ch:lec13}. (Any $A$ with
$\rho(A)<1<\norm{A}$, e.g.\ a scaled Jordan block, works equally well.)
\end{solution}

\begin{solution}[Exercise~\ref{ex:lec07-7}]
By \cref{prop:lec07-is} the sample count to reach additive error $\varepsilon$ is
$R=\bigo(\norm{c}_1^4\norm{O}^2/\varepsilon^2)$. With $\norm{c}_1=100$,
$\norm{O}\le1$, $\varepsilon=10^{-2}$,
\[
  R=\bigo\!\left(\frac{100^4\cdot1}{(10^{-2})^2}\right)
   =\bigo\!\left(\frac{10^{8}}{10^{-4}}\right)=\bigo(10^{12}),
\]
a very large number of shallow runs. By contrast, the coherent-LCU repeat count to
see the flag is $\bigo(\norm{c}_1/\norm{u})$; with $\norm{u}=\Theta(\norm{c}_1)$ this
is $\bigo(1)$ per amplification round (and oblivious amplitude amplification makes the
flag deterministic in $\bigo(\norm{c}_1/\norm{u})=\bigo(1)$ rounds), after which the
observable is estimated to error $\varepsilon$ with $\bigo(1/\varepsilon^2)=\bigo(10^4)$
measurement shots---eight orders of magnitude fewer runs than the hybrid route, though
each run is a deeper circuit carrying all $K=100$ terms and an ancilla register.

Choice: (i) on a shallow near-term device where the deep multiplexed $\mathrm{SEL}$
and clean ancillas are unavailable, the hybrid estimator is preferable despite the
large $R$, because each circuit is trivially shallow and needs no coherent control
register; (ii) if the normalized state $u/\norm{u}$ must feed a downstream quantum
subroutine, only the coherent LCU produces it---the hybrid estimator returns just the
scalar $u\dg Ou$---so coherent LCU is mandatory there.
\end{solution}
